\bookchapter{Capítulo 68: O caminho avançado — permissões, notificações e arquitetura moderna}

\chapterintro{Conseguir fazer um App é só o começo. Este capítulo fala das etapas que levam de “saber fazer” a “fazer bem”: como pedir permissões com elegância, como enviar notificações sem incomodar, como elevar a arquitetura mais um nível e — em 2026 — como programar iOS em parceria com a IA. Cada ponto é algo que você vai encontrar de imediato em projetos reais.}

\booksection{Permissões: primeiro bata na porta, depois entre}

Câmera, localização, contatos, notificações... nessas “zonas sensíveis”, o iOS exige que o App peça permissão ao usuário com um pop-up no momento do uso. A regra tem três passos: “declarar + pedir + degradar com elegância”:

\begin{lstlisting}[language=Swift]
// 1. Declare o texto de uso no Info.plist (sem isso, crash imediato!)
// Privacy - Location When In Use Usage Description: "Usado para consultar o tempo na sua cidade"

// 2. Pedir na hora do uso
import CoreLocation

@Observable
class LocationModel {
    private let manager = CLLocationManager()

    func request() {
        switch manager.authorizationStatus {
        case .notDetermined: manager.requestWhenInUseAuthorization()
        case .denied, .restricted: /* orientar o usuário às Configurações */
        default: /* já autorizado, seguir em frente */
        }
    }
}

// 3. Após a recusa, degrace com elegância: aviso + abrir Configurações
// UIApplication.shared.open(URL(string: UIApplication.openSettingsURLString)!)
\end{lstlisting}

\begin{tcbwarn}{Atenção: a permissão no iOS vale "uma vez só"}
O pop-up de permissão do iOS aparece apenas uma vez — depois que o usuário recusa, ele nunca mais aparece. Por isso: escolha o momento “em que o usuário está prestes a usar” (e não empilhe pop-ups logo na abertura) e escreva um texto explicando a finalidade. A orientação na tela após uma recusa (“ative a localização em Configurações”) é item obrigatório — não deixe o usuário preso em um beco sem saída.

\end{tcbwarn}

\booksection{Notificações: dizer algo útil na hora certa}

A notificação é o canal do App para alcançar o usuário — e também o recurso mais fácil de ser “bloqueado”. No iOS usa-se o framework UserNotifications, com três regras de ouro: enviar só o que é útil, permitir que o usuário desligue e não importunar de madrugada. Uso básico:

\begin{lstlisting}[language=Swift]
import UserNotifications

// 1. Pedir permissão de notificação (diálogo único)
UNUserNotificationCenter.current()
    .requestAuthorization(options: [.alert, .sound, .badge]) { granted, _ in
        print(granted ? "autorizado" : "negado")
    }

// 2. Agendar um "lembrete de despesas diário às 20:00"
func scheduleDailyReminder() {
    let content = UNMutableNotificationContent()
    content.title = "Lembrete de despesas"
    content.body = "Você ainda não registrou nada hoje"
    content.sound = .default

    var comps = DateComponents()
    comps.hour = 20
    let trigger = UNCalendarNotificationTrigger(dateMatching: comps, repeats: true)

    let request = UNNotificationRequest(
        identifier: "daily-reminder",
        content: content,
        trigger: trigger
    )
    UNUserNotificationCenter.current().add(request)
}
\end{lstlisting}

Atenção: a permissão de notificações no iOS também aparece “só uma vez”, e o usuário pode desativá-la a qualquer momento nas Configurações — por isso o valor percebido do conteúdo é essencial: um único excesso de envio pode fazer você perder o usuário para sempre.

\booksection{Arquitetura avançada: aprofundando o MVVM e a testabilidade}

O Capítulo 65 já usou o MVVM; aqui vão três pontos avançados para complementar:

\begin{itemize}
  \item Abstração por protocolo: faça o ViewModel depender de um protocol em vez de uma classe concreta — nos testes dá para injetar uma "fonte de dados falsa"; é a base da testabilidade;
  \item Responsabilidade única: um ViewModel cuida de uma coisa só (o de anotações não cuida do tempo); quando as funções aumentam, separe;
  \item Não exagere na engenharia: em um projeto pequeno de uma pessoa só, "classe @Observable + abstração por protocolo" já basta. Arquiteturas pesadas como VIPER e Clean Architecture ficam para quando a equipe e a complexidade baterem à porta — a ferramenta serve para atingir objetivos, não para exibir.
\end{itemize}

\begin{tcbinfo}{Dica: os "dividendos do ecossistema" do SwiftUI}
O ecossistema iOS ainda tem dois recursos avançados exclusivos: os widgets (WidgetKit) permitem ver informações sem abrir o App (coloque o “gasto do mês” do livro-caixa na tela inicial) e os atalhos (App Intents) permitem acionar suas funções pela Siri ou por automações. Ambos têm frameworks oficiais, custo de entrada baixo e um ganho enorme na experiência do usuário — no caminho avançado, vale desbloqueá-los primeiro.

\end{tcbinfo}

\booksection{A novidade de 2026: desenvolvimento iOS assistido por IA}

Em 2026, a IA já está profundamente integrada ao fluxo de desenvolvimento da Apple. Se você quer ganhar velocidade, esta combinação é bem madura:

\begin{center}
\begin{tabularx}{\linewidth}{|X|X|X|}
\hline
\textbf{Etapa} & \textbf{Como a IA ajuda} & \textbf{Sua responsabilidade} \tabularnewline
\hline
Escrever código & Autocompletar integrado ao Xcode + geração por plugins de IA de terceiros & Descrever o requisito com clareza e revisar o resultado \tabularnewline
Pesquisar & Entregar a documentação oficial e as mensagens de erro para a IA interpretar & Conferir as conclusões importantes na documentação oficial \tabularnewline
Escrever testes & Pedir à IA para gerar casos de teste XCTest nas três categorias "normal/limite/exceção" & Confirmar que os casos fazem sentido para o negócio \tabularnewline
Corrigir bugs & Colar o log de crash + o código relacionado e deixar a IA localizar e corrigir & Entender a mudança e rodar os testes para validar \tabularnewline
\hline
\end{tabularx}
\end{center}

\begin{center}\small\color{gray}Os quatro cenários do desenvolvimento iOS assistido por IA\end{center}

A metodologia concreta — como dar instruções à IA, como validar os resultados e como evitar armadilhas — está nos 11 capítulos completos da Parte 7 deste livro, “Prática de Programação com IA” (Capítulos 37-47); a leitura em paralelo é fortemente recomendada: deixe a sintaxe do Swift com esta parte e os segredos da colaboração com IA com a Parte 7.

\begin{tcbpractice}{Mãos à obra: adicione o "kit de três itens" ao livro-caixa}
\booksection{Essenciais da concorrência em Swift: @MainActor e Task}

O sistema de concorrência do Swift 6 tem duas regras que todo iniciante precisa entender:

\begin{itemize}
  \item A interface só pode ser atualizada na thread principal: as views do SwiftUI são @MainActor por padrão; nas respostas de rede, volte à thread principal antes de mudar o estado (com await MainActor.run \{ ... \} ou marcando o ViewModel como @MainActor);
  \item Gerencie as tarefas assíncronas com Task: o modificador .task \{ \} começa sozinho quando a view aparece e é cancelado automaticamente quando ela desaparece — essa é a forma padrão de iniciar trabalho assíncrono no SwiftUI; nada de espalhar Tasks soltos no body.
\end{itemize}

\begin{lstlisting}[language=Swift]
struct WeatherScreen: View {
    @State private var vm = WeatherViewModel()

    var body: some View {
        Group {
            switch vm.state {
            case .loading: ProgressView()
            case .success(let d): WeatherContent(data: d)
            case .error(let m): Text(m)
            }
        }
        .task {                       // começa ao aparecer, cancela ao sumir
            await vm.search(city: "Pequim")
        }
    }
}
\end{lstlisting}

Por trás dessas duas regras está o design de “segurança contra data races” do Swift 6 — o compilador acusa o erro na hora em que você escreve errado, em vez de esperar o App travar para descobrir. O erro do compilador não é implicância: é ele tomando a bala por você.

(1) Adicione ao livro-caixa um “lembrete de anotar os gastos às 20:00 todos os dias” (permissão de notificação + notificação local); (2) transforme a fonte de dados do ViewModel em abstração por protocolo e escreva um teste unitário com dados falsos injetados; (3) crie um widget de tela inicial com o “gasto do mês”. Com o kit completo, seu projeto ganha cara de “profissional de verdade”.

\end{tcbpractice}

\begin{tcbsummary}{Resumo do capítulo}
\begin{itemize}
  \item Permissões em três passos: declarar no Info.plist (se faltar, o App trava na hora), pedir no momento do uso e, se recusado, orientar a ir às Configurações
  \item O pop-up de permissão do iOS só aparece uma vez: escolha o momento "em que o usuário está prestes a usar" e explique com clareza a finalidade
  \item Três regras das notificações: só enviar o que é útil, poder desligar e não importunar; framework UserNotifications
  \item MVVM avançado: abstração por protocolo (testável), responsabilidade única e sem engenharia em excesso
  \item Dividendos do ecossistema: widgets com WidgetKit + atalhos com App Intents
  \item Quatro cenários com IA: escrever código, pesquisar, escrever testes e corrigir bugs; a metodologia está na Parte 7 (Capítulos 37-47)
\end{itemize}

\end{tcbsummary}
